\section{確率的プログラミングの応用3;状態空間モデル}
\subsection{授業内容}
\subsubsection{科目の中でのこのコマの位置づけ}
前時に時系列的なデータを導入したが，時系列的な性質を無視したモデリングになっていた。時系列的な分析方法は，心理学においてもウェアラブル端末の利用やSNSなど公的なデータを分析することなどにも利用できるため，非常に有用なものになりうる。しかしデータの特徴として非独立性の問題，周期性やトレンドの存在などがあり，周波数解析をおこなったり多次元の行列分解などが必要である。中でも状態空間モデルは比較的シンプルであり，とくにベイジアンアプローチで実装が容易になったと言えるだろう。

ここでは状態空間モデルの基本的な考え方を導入し，モデリングについて解説する。
ここで状態と観測の分離を行い，とくに観測が行われていない点があっても分析できること，観測が行われていない点をパラメータとして保管することに言及する。
観測が行われていない点が保管できるのであれば，未来の時点についても予測が可能になるということである。ホワイトノイズモデルでそれを行うと，確信区間が広がっていくことが観測される。
そこでトレンドを入れたモデルにすることで，さらに予測の形を変えられることを学ぶ。

そのほかにも季節項など，時系列特有の情報を組み込めることや，二次元に展開することで空間データの分析にも応用できることに言及する。
\subsubsection{コマ主題細目}
\begin{description}
	\item[時系列データの特徴] 時系列的なデータがどのように得られ，どのようなシーンで利用可能であるかを概観する。ここで時系列データはサンプルの独立性が満たされていないという問題があるため単純な線形回帰は不適切であること，また周期性やトレンド，介入効果が出てくるまでの期間など独自に考えなければならないことがいろいろ含まれている。これまで研究されてきた領域や研究方法について概観する。
	\item[状態空間モデル] さまざまな分析方法がこれまで考えてこられているが，状態空間モデルはその中でも比較的簡単な数理的構造を持ち，またベイジアンモデリングを利用することでかつての分析モデルが必要としたスムージングなどを，特段意識することなく分析できる。状態と観測というモデルの基本構造を提示し，これらがどのように実装可能かをみる。
	      \begin{flushright}
		      $\to$\citet{Matsuura201610}のPp.229-235,\citet{Baba2019}の第5部
	      \end{flushright}
	\item[欠損値の補間] 観測時点には欠損が含まれることもあり，これをパラメータとして推定・補間することができる。またこれが可能であるということは，未来の時点を欠損値として考えれば予測ができることにもなる。プログラミングの工夫により，欠損を補間するようなコードの書き方を学ぶ。また単純なホワイトノイズモデルであればあまり予測として意味がないが，トレンドを考えることで時系列的な影響について考えることができる。
	\item[状態空間モデルの展開] 状態空間モデルは，説明変数を加えた回帰モデルに応用したり，周期性をモデリングすることなども可能である。さらに時系列は一次元的であるが，二次元にも広げると空間分析にも利用が可能であることに言及する。これからの心理学は，時系列や空間など状況変数をより積極的に取り組んだモデルも利用するようになるだろう。
\end{description}
\subsubsection{キーワード}
\begin{itemize}
	\item 状態空間モデル
	\item トレンド
	\item 補間
\end{itemize}
\subsection{授業情報}
\paragraph{コマの展開方法}
講義/遠隔可/演習
\subsubsection{予習・復習課題}
\paragraph{予習} 時系列データを，時間を独立変数とした回帰分析にすることでどういった問題があるのかについて，回帰分析や確率モデルの前提などを考えて振り返っておくことが良い予習になるだろう。
\paragraph{復習} 身の回りの身近ところからでもデータを取ることができるのが，時系列データのおもしろいところでもあるので，応用可能なデータを探して分析してみると良い。可能であれば今日からでも，時系列的なデータを取り始めると，長期的に見て非常に興味深い分析ができるようになるだろう。